% Paper 2: Context Distraction Effect in Small Language Models
% Target: EMNLP / ACL Findings / NeurIPS Workshop

\documentclass[11pt]{article}
\usepackage[utf8]{inputenc}
\usepackage{amsmath,amssymb}
\usepackage{graphicx}
\usepackage{booktabs}
\usepackage{hyperref}

\title{Context Distraction: Why Retrieval-Augmented Generation\\
Degrades Small Language Model Performance\\on Structured Code Tasks}

\author{William Chen\\
National Taiwan University}

\date{}

\begin{document}
\maketitle

\begin{abstract}
We present evidence that retrieval-augmented generation (RAG) degrades
the performance of small language models ($<$10B parameters) on
structured code generation tasks. In HLS pragma optimization,
providing few-shot examples reduces type match accuracy from 90.0\%
to 73.4\% and Pragma Fidelity Score from 0.82 to 0.75 compared to
zero-shot prompting. We term this phenomenon \emph{Context Distraction}:
small models with limited attention capacity attend to surface patterns
in retrieved examples rather than reasoning about the target code's
structure. Multi-model ensembles and prompt chaining similarly fail
to improve over simple zero-shot prompts. Our findings suggest that
for small local models on domain-specific tasks, less context is more.
\end{abstract}

\section{Introduction}

Retrieval-Augmented Generation (RAG) has become the default approach
for improving LLM performance on domain-specific tasks. The assumption
is straightforward: providing relevant examples helps the model produce
better outputs. This assumption holds for large models (100B+
parameters) with extensive context windows.

We challenge this assumption for small models ($<$10B parameters)
deployed in constrained environments. In our study of HLS pragma
optimization---a structured code generation task where models must
suggest hardware optimization directives---we find that:

\begin{enumerate}
    \item \textbf{RAG hurts}: Few-shot examples reduce accuracy from
    90.0\% to 73.4\% (Table~\ref{tab:main}).
    \item \textbf{Chaining hurts Gemma4}: Multi-step prompting drops
    PFS from 0.883 to 0.813 for MoE models (5-kernel subset).
    \item \textbf{Multi-model fusion fails}: Combining outputs from
    multiple models performs worse than the best single model.
    \item \textbf{Zero-shot is optimal}: Simple, focused prompts
    consistently outperform all augmentation strategies.
\end{enumerate}

\section{Experimental Setup}

We evaluate on HLS pragma prediction using ForgeHLS (184K designs).
Models: Gemma4 (4.5B MoE), Qwen 2.5 Coder (7B dense), DeepSeek-R1
(14B dense). All run locally via Ollama.

\subsection{Prompting Strategies}

\begin{itemize}
    \item \textbf{A1 (Zero-shot)}: Direct optimization request with
    no examples.
    \item \textbf{A3 (Rule-based)}: Optimization request plus HLS
    best practices (5 rules).
    \item \textbf{A4 (Few-shot RAG)}: 3 retrieved similar kernels
    with optimal pragmas as examples.
    \item \textbf{A5 (Aggressive)}: Zero-shot with strong optimization
    directive.
\end{itemize}

\section{Results}

\begin{table}[t]
\centering
\caption{Prompting strategy comparison on HLS pragma prediction.
10 kernels, Qwen 2.5 Coder 7B. More context consistently degrades
performance.}
\label{tab:main}
\begin{tabular}{lcccc}
\toprule
\textbf{Strategy} & \textbf{Context} & \textbf{Type Match} & \textbf{PFS} \\
\midrule
A1: Zero-shot      & None      & \textbf{90.0\%} & \textbf{0.818} \\
A5: Aggressive      & Directive & \textbf{90.0\%} & \textbf{0.820} \\
A3: HLS Rules       & 5 rules   & 86.7\%          & 0.781 \\
A4: Few-shot RAG    & 3 examples & 73.4\%          & 0.752 \\
\bottomrule
\end{tabular}
\end{table}

\begin{table}[t]
\centering
\caption{Context distraction across models and strategies.
Fair comparison: same inference count per group.}
\label{tab:models}
\begin{tabular}{lccc}
\toprule
\textbf{Strategy} & \textbf{Inferences} & \textbf{PFS} \\
\midrule
Gemma4 pass@1 (zero-shot)   & 1 & \textbf{0.883}$^\dagger$ \\
Gemma4 chaining              & 2 & 0.813 \\
Qwen$\to$Gemma4 dual         & 2 & 0.697 \\
Gemma4$\to$Qwen dual         & 2 & 0.521 \\
Qwen pass@1 (zero-shot)      & 1 & 0.633 \\
Qwen chaining                & 2 & 0.527 \\
\bottomrule
\end{tabular}
\end{table}

\subsection{Key Observation}

Every strategy that adds context---whether retrieved examples (RAG),
explicit rules, chaining steps, or multi-model collaboration---performs
\emph{worse} than zero-shot. The degradation is consistent across
both dense (Qwen 7B) and MoE (Gemma4) architectures.

\section{Analysis: Why Does Context Hurt?}

\subsection{Attention Competition}

Small models have limited attention capacity. When examples are added
to the prompt, the model must divide attention between:
(a) understanding the target kernel's structure, and
(b) pattern-matching against the examples.

For 7B models, (b) wins---the model copies surface patterns from
examples rather than reasoning about the target. This is
\emph{mimicry over reasoning}.

\subsection{Oracle Experiments: Ruling Out Implementation Bias}

A natural objection is that our negative results stem from poor
implementation (bad retrieval, inaccurate classification) rather than
fundamental model limitations. We address this with oracle experiments
on 20 kernels where we provide \emph{perfect} inputs:

\begin{table}[t]
\centering
\caption{Oracle experiments. Even with perfect routing and optimal
examples, additional context hurts. $^*p < 0.05$}
\label{tab:oracle}
\begin{tabular}{lcc}
\toprule
\textbf{Method} & \textbf{PFS} & \textbf{vs Baseline} \\
\midrule
Baseline (generic prompt)        & \textbf{0.696} & --- \\
Oracle Routing (manual classify) & 0.490          & $-29.6\%^*$ \\
Oracle RAG (best 1-shot)         & 0.614          & $-11.8\%$ \\
\bottomrule
\end{tabular}
\end{table}

\textbf{Even with 100\% correct classification}, domain-specific skills
are \emph{significantly} worse than generic prompts ($t=-2.78, p<0.01$).
This eliminates the ``bad implementation'' hypothesis: Context Distraction
is an inherent property of small model attention capacity, not an artifact
of our experimental setup.

\subsection{Cross-Model Pragma Synthesis Failure}

We tested merging pragma suggestions from multiple models
(CMPS): taking PIPELINE/UNROLL from Qwen and PARTITION from Gemma4.
This rule-based merge achieved PFS=0.690, \emph{worse} than either
model alone. The failure stems from stochastic output variation:
a model's strength in one pragma type is not consistent across runs.

\section{Implications}

\begin{enumerate}
    \item \textbf{For practitioners}: When deploying small models
    ($<$10B) for structured tasks, use simple zero-shot prompts.
    RAG and chaining are counterproductive.
    \item \textbf{For researchers}: The ``more context is better''
    assumption should be tested for each model size class. Context
    distraction may explain why RAG gains diminish for smaller models.
    \item \textbf{For air-gapped deployment}: Zero-shot being optimal
    is good news---no retrieval infrastructure needed for local models.
\end{enumerate}

\section{Related Work}

\textbf{Retrieval-Augmented Generation.}
RAG~\cite{lewis2020rag} has become standard for knowledge-intensive
NLP tasks, demonstrating consistent improvements for models with
100B+ parameters. However, its effectiveness on small models ($<$10B)
for structured code generation remains underexplored. Our work provides
the first systematic evidence that RAG can be counterproductive for
small models on domain-specific tasks.

\textbf{In-Context Learning Sensitivity.}
Recent work shows that ICL performance depends heavily on example
selection, ordering, and formatting~\cite{wei2022cot}. We extend
this finding by showing that even with \emph{optimal} example selection
(oracle RAG), small models still degrade---suggesting the issue is
attention capacity, not example quality.

\textbf{LLMs for Hardware Design.}
ChipNeMo~\cite{chipnemo} fine-tunes domain-specific LLMs but targets
natural language tasks (documentation, bug summarization) rather than
code generation. HLSPilot uses GPT-4 for pragma suggestion but relies
on cloud APIs. Our work focuses on local small models where context
efficiency is critical due to limited attention capacity.

\textbf{Attention Capacity in Small Models.}
The ``lost in the middle'' phenomenon~\cite{liu2024lost} shows that
even large models struggle to use information in long contexts. We
demonstrate a more severe version for small models: \emph{any}
additional context---even highly relevant domain knowledge---degrades
performance on structured generation tasks.

\section{Conclusion}

We identify Context Distraction as a systematic failure mode of
small language models: providing more context degrades performance
on structured code tasks. Zero-shot prompting with post-processing
is the optimal strategy for local models under 10B parameters.

\bibliographystyle{plain}
\begin{thebibliography}{9}
\bibitem{lewis2020rag} P. Lewis et al., ``Retrieval-Augmented Generation for Knowledge-Intensive NLP Tasks,'' NeurIPS, 2020.
\bibitem{wei2022cot} J. Wei et al., ``Chain-of-Thought Prompting Elicits Reasoning in Large Language Models,'' NeurIPS, 2022.
\bibitem{chipnemo} M. Liu et al., ``ChipNeMo: Domain-Adapted LLMs for Chip Design,'' DAC, 2024.
\bibitem{liu2024lost} N. Liu et al., ``Lost in the Middle: How Language Models Use Long Contexts,'' TACL, 2024.
\end{thebibliography}

\end{document}
